\section{Relações de ordem}
Nesta seção, introduziremos as noções básicas de pré-ordem, ordem parcial e ordem total, bem como suas versões estritas e conceitos associados.

\begin{definition}[Pré-ordem]\index{pré-ordem}
    Uma \emph{pré-ordem} é um par $(A, \leq)$ onde $A$ é um conjunto e $\leq$ é uma relação binária em $A$ tal que para todo $x, y, z \in A$:
    \begin{itemize}
        \item $x \leq x$ (reflexividade)
        \item Se $x \leq y$ e $y \leq z$, então $x \leq z$ (transitividade).
    \end{itemize}
\end{definition}

\begin{definition}[Ordem parcial]\index{ordem parcial}
    Uma \emph{ordem parcial} é um par $(A, \leq)$ onde $A$ é um conjunto e $\leq$ é uma pré-ordem e, além disso, satisfaz a seguinte condição.
    \begin{itemize}
        \item Se $x \leq y$ e $y \leq x$, então $x = y$ (antissimetria).
    \end{itemize}
\end{definition}

\begin{definition}[Ordem total]\index{ordem total}\index{ordem linear}
    Uma \emph{ordem total}, ou \emph{ordem linear} é um par $(A, \leq)$ onde $A$ é um conjunto e $\leq$ é uma ordem parcial e, além disso, satisfaz a seguinte condição.
    \begin{itemize}
        \item Para todo $x, y \in A$, ou $x \leq y$ ou $y \leq x$ (totalidade).
    \end{itemize}
\end{definition}

Há também versões estritas.

\begin{definition}
    Uma \emph{ordem parcial estrita} é um par $(P, <)$ onde $P$ é um conjunto e $<$ é uma relação binária em $P$ tal que para todo $x, y, z \in P$, valem as seguintes condições.
    \begin{itemize}
        \item Não existe $x$ tal que $x < x$ (irreflexividade).
        \item Se $x < y$ e $y < z$, então $x < z$ (transitividade).
    \end{itemize}
\end{definition}

\begin{definition}
    Uma \emph{ordem total estrita} é um par $(P, <)$ onde $P$ é um conjunto e $<$ é uma ordem parcial estrita e, além disso, satisfaz a seguinte condição.
    \begin{itemize}
        \item Para todo $x, y \in P$, ou $x < y$ ou $y
         < x$ ou $x = y$ (totalidade).
    \end{itemize}
\end{definition}

As versões estritas e não estritas estão em correspondência natural.

\begin{definition}[Ordem parcial estrita induzida]
    Seja $(P, \leq)$ uma ordem parcial.
    A ordem parcial estrita associada a $(P, \leq)$ é o par $(P, <)$ em que $<$ é a relação binária em $P$ definida por $x < y$ se e somente se $x \leq y$ e $x \neq y$.
\end{definition}

\begin{definition}[Ordem parcial não estrita induzida]
    Seja $(P, <)$ uma ordem parcial estrita.
    A ordem parcial associada a $(P, <)$ é o par $(P, \leq)$ em que $\leq$ é a relação binária em $P$ definida por $x \leq y$ se e somente se $x < y$ ou $x = y$.
\end{definition}

As seguintes duas proposições são deixadas como exercício ao leitor (Exercício~\ref{ex:ordens-duais}).

\begin{proposition}\label{prop:ordens-duais}
    Seja $(P, \leq)$ uma ordem parcial.
    \begin{enumerate}[label=(\alph*)]
        \item A ordem parcial estrita induzida por $(P, \leq)$ é de fato uma ordem parcial estrita.
        \item Seja $(P, <)$ a ordem parcial estrita induzida por $(P, \leq)$.
        Então a ordem parcial não estrita induzida por $(P, <)$ é $(P, \leq)$.
        \item $(A, \leq)$ é uma ordem total se e somente se $(A, <)$ é uma ordem total estrita.
    \end{enumerate}
\end{proposition}

Dualmente, temos a seguinte proposição.

\begin{proposition}\label{prop:ordens-duais-2}
    Seja $(A, <)$ uma ordem parcial estrita.
    \begin{enumerate}[label=(\alph*)]
        \item A ordem parcial induzida por $(A, <)$ é de fato uma ordem parcial.
        \item Seja $(A, \leq)$ a ordem parcial induzida por $(A, <)$.
        Então a ordem parcial estrita induzida por $(A, \leq)$ é $(A, <)$.
        \item $(A, <)$ é uma ordem total estrita se e somente se $(A, \leq)$ é uma ordem total.
    \end{enumerate}
\end{proposition}

Por isso, ao se estudar ordens parciais, não há perda em se restringir ao estudo de ordens estritas ou não estritas.
Além disso, ao longo deste livro, ao estabelecer que $\leq$ é uma ordem parcial, iremos, sem aviso, utilizar a ordem estrita $<$ associada a $\leq$ e vice-versa.
Outros símbolos que poderão ser usados para denotar ordens parciais fazendo jus a esta dualidade são $\preceq$ e $\prec$, $\sqsubseteq$ e $\sqsubset$, $\leqslant$ e $\lneqq$, entre outros.

Abaixo, definiremos os conceitos de máximo, mínimo, maior elemento e menor elemento de um conjunto em uma pré-ordem.
Para ordens parciais estritas, tais conceitos são definidos a partir da ordem parcial não estrita associada.

\begin{definition}
    Seja $(P, \leq)$ uma pré-ordem e $A \subseteq P$. Dizemos que:
    \begin{itemize}
        \item $x$ é uma \emph{cota superior} de $A$ se para todo $a \in A$, $a \leq x$;
        \item $x$ é uma \emph{cota inferior} de $A$ se para todo $a \in A$, $x \leq a$;
        \item $x$ é um \emph{máximo} de $A$ se é uma cota superior de $A$ e $x \in A$.
        \item $x$ é um \emph{mínimo} de $A$ se é uma cota inferior de $A$ e $x \in A$.
        \item $x$ é \emph{maximal} em $A$ se para todo $a \in A$, tal que $x \leq a$, segue que $a\leq x$.
        \item $x$ é \emph{minimal} em $A$ se para todo $a \in A$, tal que $a \leq x$, segue que $x \leq a$.
        \item $x$ é um \emph{supremo} de $A$ se é a menor cota superior de $A$, ou seja, $x$ é uma cota superior de $A$ e para toda cota superior $y$ de $A$, segue que $x \leq y$.
        \item $x$ é um \emph{ínfimo} de $A$ se é a maior cota inferior de $A$, ou seja, $x$ é uma cota inferior de $A$ e para toda cota inferior $y$ de $A$, segue que $y \leq x$.
    \end{itemize}
\end{definition}
\subsection*{Exercícios}
\begin{exercise}\label{ex:ordens-duais}
    Prove as proposições \ref{prop:ordens-duais} e \ref{prop:ordens-duais-2}.
\end{exercise}